% Section 5.3: Deep-Dive Case Studies Across E1-E4

\subsection{Deep-Dive Case Studies}

\subsubsection{Case Study 1: Hidden Single-Scan Delay Masked by 100\% Simulation Pass (\code{PUB011})}
\code{PUB011} specifies a three-station transfer trolley toggling between manual and automatic operation. Cross-condition inspection contrasts both runtime-vs.-formal coverage and per-run synthesis variance:
\begin{itemize}
  \item \textbf{E3 Passes 8/8 Fixed Tests Yet Fails 8/14 NuSMV Properties}: DeepSeek under STAR (\pathcode{results/raw/four_group_20261114/star_deepseek_v3/PUB011/generation/compile_00/FB_Task.st}) achieves \textbf{8/8 (100\%)} pass on real CODESYS~1\,ms simulation tests. During Stage~3 verification, however, \code{NuSMV} proves only 6/14 LTLs and builds deterministic counterexamples against 8 transition obligations (incl.\ \code{SMV_4}, \code{SMV_8}, \code{SMV_9}) due to an intermediate \code{E_Phase.MANUAL} state inserted before the dispatcher:
\begin{lstlisting}[caption={Relevant E3 (\code{PUB011}) mode-switch logic introducing a one-scan startup delay.},float=htbp]
IF NOT Mode THEN
    phase := E_Phase.MANUAL;
ELSIF phase = E_Phase.MANUAL THEN
    phase := E_Phase.WAIT_A;
ELSE
    CASE phase OF ...
\end{lstlisting}
  On the exact rising-edge scan transitioning from \code{Mode = FALSE} to \code{TRUE} with \code{SensorA} and \code{ButtonA} asserted, one scan is spent resetting \code{phase} from \code{MANUAL} to \code{WAIT_A}, delaying \code{MotorForward} by one scan cycle---a corner case never triggered by the 8 finite test sequences.
  \item \textbf{E1, E2, and E4 Preserve Same-Scan Entry Semantics}: By contrast, \textbf{E4} (\pathcode{results/raw/four_group_20261114/star_terra/PUB011/generation/compile_00/FB_Task.st}, maintaining \code{E_TaskPhase.WAIT_A} during \code{NOT Mode}) as well as both \code{Codex CLI} baselines (\textbf{E1} and \textbf{E2}) stay directly in their \code{WAIT_A} phase while in manual mode, allowing immediate same-scan transition and simultaneously passing \textbf{8/8 simulation cases and 14/14 formal LTLs}.
\end{itemize}

\subsubsection{Case Study 2: Universal Transient Latching and Priority Arbitrations (\code{PUB019})}
On the 18-input, 6-output multi-mode robotic arm controller \code{PUB019} (Hard), \textbf{all four conditions (\textbf{E1}--\textbf{E4}) converge on the exact same residual flaw}: every group scores \textbf{16/17} on fixed cases, failing exclusively at scan~1 of \code{PUB019.T6} (\code{MoveUp} expected \code{TRUE}, measured \code{FALSE}), and every group terminates with Stage~3 \code{COUNTEREXAMPLE} (\textbf{E1}: 4/8, \textbf{E2}: 3/8, \textbf{E3}: 5/10, \textbf{E4}: 6/10 properties satisfied). Two common failure mechanisms explain this convergence:
\begin{enumerate}
  \item \textbf{Over-Suppression on Entry From Idle}: At scan~0 of \code{PUB019.T6}, no mode is active; scan~1 raises \code{ManualMode} alongside \code{UpButton}. Across all four conditions, models detect \code{modeChanged} against quiescent idle state and blank actuator outputs for one scan, rather than gating only active-to-active mode switches.
  \item \textbf{Conservative Mutual Lockout vs.\ Manual Preemption}: When \code{ManualMode} and an auto-mode selector activate concurrently, all four candidates apply \code{modeCount > 1} to freeze all axes safely, whereas the specification's observer defines explicit manual override precedence, triggering immediate scan-0 formal counterexamples.
\end{enumerate}

\subsubsection{Case Study 3: Formal Catch on Bare Codex (\code{PUB020}), Model Divergence (\code{PUB016}), and Indexing Nuance (\code{PUB014})}
\begin{itemize}
  \item \textbf{\code{PUB020} (Multi-Station Production Line)}: Both \code{gpt-6.1-terra} conditions (\textbf{E2} and \textbf{E4}) pass \textbf{11/11 fixed cases and Stage~3 \code{PASS}}. Conversely, both DeepSeek runs draw Stage~3 counterexamples: critically, \textbf{E1} passes \textbf{all 11/11 finite simulation cases} yet satisfies only 3/11 \code{NuSMV} properties (exposing mutual-exclusion violations in \code{REJECT} reachable states, complementing \textbf{E3}/\code{PUB011}), while \textbf{E3} drops late progression in \code{PUB020.T9} (\code{10/11} cases, \code{3/10} formal).
  \item \textbf{\code{PUB016} (Analog Ramp/Step Control)}: Inverting overall model rank order, both DeepSeek groups (\textbf{E1} and \textbf{E3}) pass \textbf{8/8} tests, whereas both GPT groups (\textbf{E2} and \textbf{E4}) fail identically on \code{PUB016.T4} (\textbf{7/8}).
  \item \textbf{\code{PUB014} (Dotted-Decimal IPv4 Parser)}: Both DeepSeek runs (\textbf{E1} and \textbf{E3}) compile on the first attempt yet score \textbf{0/8} at runtime because indexing \code{ipText[i]} defaults to 1-based \code{1..LEN(ipText)} instead of CODESYS's \textbf{0-based} \code{STRING} byte offset \code{s[i]}, reading past the terminator and clearing output. (Both GPT conditions hit compiler-lock timeouts before simulation.) Targeted updates to \pathcode{dialect_cards.py} subsequently codify this 0-based vs.\ 1-based distinction explicitly.
\end{itemize}

\subsubsection{Case Study 4: Vendor Static-Analysis Guardrails (\code{SA0027}/\code{SA0033}) and Generic Reflection (\code{PUB018})}
\begin{itemize}
  \item \textbf{Same-Root \code{SA0027} Elimination Across \textbf{E1} vs.\ \textbf{E3}}: All four compilation failures in \textbf{E1} (\code{Codex + DeepSeek}: \code{PUB003}, \code{PUB009}, \code{PUB017}, \code{SELF009}) failed on the exact same CODESYS static-analysis collision---\code{SA0027} (\code{Variable name 'error' in 'FB_Task'} \code{already used for an object in Library} \code{'caa device diagnosis, 3.5.22.0'}), compounded by \code{SA0033} on \code{PUB017}---despite multiple multi-turn repair iterations. By contrast, \textbf{E3} (\textbf{STAR} + DeepSeek) compiled \textbf{4/4 with zero \code{SA0027} failures} (three on first attempt, \code{PUB017} after one repair) via dialect cards and \code{compact_diagnostics()}, while \textbf{E2} (\code{Codex + Terra}) resolved the first three after 1--2 iterative fixes.
  \item \textbf{\code{PUB018} (Generic Circular Byte Queue)}: Binding caller-owned \code{DINT} variables and fixed arrays via \code{VAR_IN_OUT VARIANT} / \code{VAR_INFO} marks an infrastructure and abstraction frontier: GPT conditions (\textbf{E2}, \textbf{E4}) hit \code{COMPILE_FAIL} (\code{Unknown type: 'VARIANT'}), whereas DeepSeek conditions (\textbf{E1}, \textbf{E3}) compile cleanly yet incur external harness \code{TOOL_ERROR} during simulation type binding.
\end{itemize}